% =====================================================================
% Cover letter VARIANT C (lead-with-honest-correction) -- NMR in Biomedicine
% EIC: Professor John R. Griffiths
% Build: tectonic cover_letter_correction.tex
% =====================================================================
\documentclass[11pt]{letter}
\usepackage[margin=1in]{geometry}
\usepackage{microtype}
\usepackage[hidelinks]{hyperref}

\signature{Avery Karlin, BA \\ Department of Applied Physics and Applied Mathematics,
Columbia University \\ \texttt{ak5232@columbia.edu}}
\address{Avery Karlin \\ Columbia University \\ New York, NY, USA}

\begin{document}
\begin{letter}{Professor John R.\ Griffiths \\ Editor-in-Chief, \textit{NMR in Biomedicine}}
\opening{Dear Professor Griffiths,}

I am pleased to submit ``Boundary-railing of conventional NLLS fits as an
assumption-free pseudo-diffusion identifiability diagnostic in IVIM MRI'' for
consideration in \textit{NMR in Biomedicine}.

I want to be transparent about the history of this work, because the correction is
part of what makes it worth publishing. An earlier version led with a calibration
claim --- a dramatic marginal under-coverage of the pseudo-diffusion coefficient
$D^{*}$ --- and a methods review rightly objected on two grounds: that an empirical
coverage statement is only as trustworthy as the noise and forward model behind its
reference truth, and that the dramatic figure depended on choices that were not
documented. We took both objections seriously and rebuilt the contribution.

Two things changed. First, an independent clean-room reimplementation showed that
the dramatic \emph{marginal} headline does not reproduce; it is recoverable only
under an undocumented, overconfident convention for the standard deviation of a
railed (unidentified) voxel. We have dropped that headline. The honest, reproducible
statement is \emph{conditional}: under a stated Cram\'er--Rao convention, symmetric
Gaussian intervals under-cover $D^{*}$ specifically in the high-$D^{*}$ tercile
(central-95\% coverage 0.63) while remaining near-nominal elsewhere; a shape-aware
quantile interval restores the marginal coverage but not the conditional residual,
and an amortized neural posterior dominates the railed baseline. The convention that
manufactured the original severity is now documented in full, alongside all dataset
identifiers, fitting and training detail, and the Cram\'er--Rao approximation.

Second --- and this is the heart of the resubmission --- we promote to the primary
claim a result that \emph{cannot} be overextended because it assumes nothing.
A conventional box-constrained least-squares IVIM fit drives $D^{*}$ onto a
parameter bound in 54.7\% of high-SNR voxels on the open OSIPI abdomen scan (95\%
CI 52.2--57.1), reproduced by the independent reimplementation (54.2\%) and
replicated across the full ROI (47.8\%) and an independent liver DWI cohort
(43.7--73.4\%). This boundary-railing is read directly off the optimiser, with no
ground truth and no noise-model trust, and so it answers the reviewer's central
objection rather than defending against it.

The result is a stronger, narrower, and fully reproducible paper: a primary claim
immune to the overextension critique, and a calibration result kept only within the
scope its evidence supports. Every number traces to a single seeded reproduction
that regenerates with one command on openly licensed data (OSIPI TF2.4, Zenodo
14605039; TCGA-LIHC via The Cancer Imaging Archive). The work fits both principal
categories of \textit{NMR in Biomedicine} and is complementary to learned-uncertainty
frameworks for IVIM such as Casali et al.\ (2026, this journal), who themselves
report residual $D^{*}$ overconfidence --- the learned echo of the same
identifiability wall.

The manuscript is original, has not been published previously, and is not under
consideration elsewhere. The author declares no competing interests, takes full
responsibility for the integrity of the work, and includes a declaration of
AI-assisted writing. I would be glad to suggest reviewers and to answer any
questions.

\closing{Sincerely,}
\end{letter}
\end{document}
